% Clase 1 — Redes Bayesianas
% Lectura: AIMA Cap. 13
\section{Redes Bayesianas}

\begin{definition}[Red Bayesiana]
Una red bayesiana es un grafo acíclico dirigido (DAG) donde:
\begin{itemize}
    \item Cada nodo representa una variable aleatoria.
    \item Los arcos representan dependencias directas.
    \item Cada nodo $X_i$ tiene una CPT: $P(X_i \mid \text{Parents}(X_i))$.
\end{itemize}
\end{definition}

\subsection{Semántica}

La distribución conjunta se factoriza como:
$$P(X_1, \ldots, X_n) = \prod_{i=1}^{n} P(X_i \mid \text{Parents}(X_i))$$

\subsection{Inferencia Exacta}

\begin{itemize}
    \item \textbf{Enumeración}: suma sobre todas las variables ocultas. Exponencial en el peor caso.
    \item \textbf{Eliminación de variables}: orden de eliminación importa para la eficiencia.
\end{itemize}

\subsection{Inferencia Aproximada}

\begin{itemize}
    \item \textbf{Muestreo directo}: muestrea según el orden topológico de la red.
    \item \textbf{Muestreo por rechazo}: descarta muestras inconsistentes con la evidencia.
    \item \textbf{Ponderación por verosimilitud}: fija la evidencia, pondera las muestras.
    \item \textbf{MCMC / Gibbs Sampling}: muestrea cada variable condicionada al resto.
\end{itemize}

\begin{example}[Red de alarma]
    Un sistema de alarma antirrobo se activa por un robo o (raramente) por un temblor. Dos vecinos, John y Mary,
    llaman si oyen la alarma, pero a veces confunden el teléfono con la alarma o no la oyen.
\end{example}

\begin{figure}[H]
\centering
\begin{tikzpicture}[
  var/.style={draw=black, thick, circle, minimum size=13mm, font=\small}
]
\node[var] (B) at (0,3)    {Robo};
\node[var] (E) at (4,3)    {Temblor};
\node[var] (A) at (2,1.2)  {Alarma};
\node[var] (J) at (0,-0.8) {John};
\node[var] (M) at (4,-0.8) {Mary};

\draw[-{Latex},thick] (B) -- (A);
\draw[-{Latex},thick] (E) -- (A);
\draw[-{Latex},thick] (A) -- (J);
\draw[-{Latex},thick] (A) -- (M);
\end{tikzpicture}
\caption{Red bayesiana Robo-Temblor-Alarma-John-Mary. Robo y Temblor son independientes a priori, pero se vuelven
\textbf{dependientes} al condicionar sobre Alarma (\textit{explaining away}); John y Mary son condicionalmente
independientes dado Alarma.}
\end{figure}

\begin{table}[H]
    \centering
    \begin{tabular}{cc}
        \toprule
        $P(\text{Robo})$ & $P(\text{Temblor})$ \\
        \midrule
        0.001 & 0.002 \\
        \bottomrule
    \end{tabular}
    \qquad
    \begin{tabular}{ccc}
        \toprule
        Robo & Temblor & $P(\text{Alarma})$ \\
        \midrule
        V & V & 0.95 \\
        V & F & 0.94 \\
        F & V & 0.29 \\
        F & F & 0.001 \\
        \bottomrule
    \end{tabular}
    \qquad
    \begin{tabular}{cc}
        \toprule
        Alarma & $P(\text{John})$ \\
        \midrule
        V & 0.90 \\
        F & 0.05 \\
        \bottomrule
    \end{tabular}
    \qquad
    \begin{tabular}{cc}
        \toprule
        Alarma & $P(\text{Mary})$ \\
        \midrule
        V & 0.70 \\
        F & 0.01 \\
        \bottomrule
    \end{tabular}
    \caption{CPTs de la red: la de Alarma tiene $2^2=4$ filas porque tiene 2 padres; John y Mary tienen 1 padre cada
    uno (Alarma), así que sus CPTs tienen 2 filas.}
\end{table}
